\scanpage{257}

\mspage{131}

\deleted{\unsure{dans} Dans $\widetilde Z\times\widetilde Z$ on a le sous-groupe diagonal, contenant $\unsure{\mathfrak{z}|\mathfrak{Z}}$, et le quotient de celui-ci par $\mathfrak{z}$ s'identifie alors au sous-groupe $Z\hookrightarrow T_G$, en d'autres termes, on a On retrouve le sous}

La suite exacte (121) se réécrit en termes de $\widetilde T$, et \unsure{reçoit|aussi|inscrit} la suite (115) comme indiqué dans le diagramme

\diagram{\resizebox{\linewidth}{!}{\begin{tikzcd}[ampersand replacement=\&,column sep=small,row sep=large]
\& 0 \arrow[r] \& \mathfrak{z} \arrow[r] \arrow[d,equal] \& \widetilde Z \arrow[r,"c",no head] \arrow[d,hook,"{\widetilde\lambda\mapsto(\widetilde\lambda,\widetilde\lambda^{-1})}"] \& Z \arrow[d,hook,"{\substack{\mathrm{incl}.\\\unsure{ca|\alpha}}}"] \& \\
(122) \& 0 \arrow[r] \& \mathfrak{z}\arrow[r,swap,"{z\mapsto(z,z^{-1})}"] \& \widetilde Z\times\widetilde Z \arrow[r,swap,"{(\widetilde\lambda,\widetilde\mu)\mapsto u_0(\widetilde\lambda)u_1(\widetilde\mu)}"] \& T_G\arrow[r] \&1
\end{tikzcd}}}

NB si $z\in\mathfrak{z}$, on a $u_0(z)=u_1(z)$

qui permet de retrouver l'inclusion canonique $Z\subset T_G$, quand $c:\widetilde Z\longrightarrow Z$ est surjective.

Ainsi, le groupe $T_G$ et l'inclusion $T_H\subset T_G$ \unsure{c.à.d} les données a)b) de (116'), s'explicitent en termes du groupe $\widetilde Z$ et \unsure{de|du} \inserted{l'hom $c:\widetilde Z\unsure{-|\longrightarrow}Z$ du} diagramme (115), qui fournit aussi l'hom. $\delta:T_G\unsure{-|\longrightarrow}Z$ \inserted{(donnée e) de (116'))} par commutativité, soit comme le composé de $\widetilde\delta|T_G:T_G\longrightarrow\widetilde Z$ et de $c:\widetilde Z\unsure{-|\longrightarrow} Z$, dans la \deleted{\unsure{suite}} description (117), soit (dans le cas de la description (122)), par commutativité du diagramme

% archive sheet number: 256
